\abstract % Структурный элемент: РЕФЕРАТ

\keywords{ИНТЕЛЛЕКТУАЛЬНЫЙ АГЕНТ, БОЛЬШИЕ ЯЗЫКОВЫЕ МОДЕЛИ, RETRIEVAL-AUGMENTED GENERATION, LEGAL-TECH, ГЕНЕРАЦИЯ ДОГОВОРОВ, ВЕКТОРНЫЙ ПОИСК}

Объектом разработки является интеллектуальный программный агент, который по произвольному словесному описанию сделки готовит проект договора в формате DOCX и сопровождающие юридические рекомендации со ссылками на статьи нормативных правовых актов.

Цель работы~--- разработать интеллектуальный программный агент, способный по произвольному словесному описанию сделки определять её гражданско-правовой тип, в диалоговом режиме собирать у пользователя недостающие для подготовки документа сведения и формировать проект договора в формате DOCX вместе с юридическими рекомендациями со ссылками на конкретные статьи нормативных правовых актов. Существующие в области правовых технологий (legal-tech) шаблонные конструкторы не дают правовой оценки описанной сделки с опорой на действующее законодательство, а лишь подставляют пользовательские значения в заранее заготовленный текст из закрытого каталога сценариев. В работе для преодоления этих ограничений применён агентный подход с опорой на проиндексированный корпус российского законодательства.

В рамках работы спроектирован и реализован агент в виде графа состояний, в узлах которого генеративные большие языковые модели (LLM) применяются совместно с методами генерации текста, дополненной поиском по внешней памяти (retrieval-augmented generation, RAG)~\cite{lewis2020rag}; внешней памятью служит подготовленный и проиндексированный в векторной базе данных корпус гражданского кодекса~РФ и профильных федеральных законов. Программная система реализована на стеке Python с использованием веб-фреймворка FastAPI, фреймворка LangGraph, сервера локального исполнения языковых моделей Ollama, векторной базы данных Qdrant, реляционной СУБД PostgreSQL и объектного хранилища MinIO.

Результатом является целостный программный продукт. Разработанный агент в режиме диалога с пользователем уточняет недостающие сведения о сделке и формирует проект договора с автоматической самопроверкой его соответствия применимым нормам, сопровождая результат юридическими рекомендациями по заключению описанной сделки. Вокруг агента реализованы веб-интерфейс, платёжный шлюз и панель администратора для управления корпусом нормативных правовых актов и тарифными планами. Эмпирическая оценка поиска по корпусу на специально подготовленном датасете из~54 эталонных описаний сделок показала метрики верхней части выдачи \texttt{hit@3}=100~\% и \texttt{precision@3}=91.4~\% при точности классификации типа сделки 100~\%; внешняя экспертная апробация подготовленных артефактов с привлечением трёх экспертов дала средние оценки качества проекта договора и юридических рекомендаций 4.42 и 4.73 по пятибалльной шкале Лайкерта при отсутствии оценок ниже уровня <<хорошо>> для всех одиннадцати рассмотренных типов сделок. Научная новизна работы~--- в одновременном сочетании в одной программной системе трёх свойств: работы с произвольной формулировкой сделки от пользователя, опоры на проиндексированный корпус российского законодательства и автоматического контроля качества подготавливаемого проекта договора; в области правовых технологий подобное сочетание ранее не зафиксировано.

Решение ориентировано на физических лиц и субъектов малого предпринимательства, заключающих типовые сделки гражданского оборота. Его применение повышает доступность первичной правовой поддержки, позволяет пользователю самостоятельно проверить, на каких нормах права основаны полученные рекомендации, и снижает риск несоблюдения требований гражданского кодекса~РФ к форме и существенным условиям договора.
